\textbf{What an SME can take from this.} On one laptop, the choice of model moved energy per request by more than an order of magnitude (1.0 to 19.1\,J across configuration means) while moving accuracy by tens of points, and the two did not move together in a simple way: the EU-origin 7B model cost as much as the most accurate model but was no more accurate than a model a fifth of its size. Measured energy, not parameter count, is what separates these options. Placement mattered too: moving the 0.5B models from the laptop GPU to four CPU cores, a stand-in for a modest on-premise server, made them slower and costlier at the same accuracy. ``Running it at the edge'' is therefore not automatically the frugal choice; it depends on the hardware the edge actually has.

\textbf{Why routing helped only modestly.} Routing pays off when a cheap model knows when it is likely to be wrong. Here the small model was often confidently wrong, so to stay near the large model's accuracy the margin cascade had to escalate 44\% of items and still pay for the small model on each of them, and its estimated saving shrank from the 79\% the small model alone offers to 29\%. The gap between the cascades and the oracle shows where the leverage is: better confidence estimation, task-aware routing features or a verifier would convert more of that headroom into savings \cite{guo2017,kadavath2022,hybridllm}. For an SME the practical reading is that routing is worth a meaningful fraction of energy, not most of it, at today's small-model calibration, and that simply using the small model may be the better trade when a few points of accuracy are affordable.

\textbf{Sovereignty as a design decision.} The scorecard's stable ranking came from the rubric, not from the data: once self-hosting is scored as the organisation's own control, every local option beats any hosted one regardless of weights, and the EU-origin model is not favoured. A different organisation could reasonably score publisher jurisdiction, training-data transparency or supply-chain origin more finely and reach a different ranking. We think this is the right way to use such an instrument: as a transparent statement of which dimensions an organisation chooses to value, with evidence attached, rather than as a single objective number. The static sensitivity rule makes the trade-off concrete: honouring a preference for an EU-origin model on sensitive items cost accuracy without saving much energy.

\textbf{What one laptop cannot show.} The continuum in this study is short: one GPU, its own CPU and a documented but unexecuted cloud tier. The figures say nothing about data-centre efficiency, batching servers or network energy, and the edge tier shares the chip with the device tier. The contribution is a method and a calibrated first data point that a larger study can extend to real edge hardware and hosted tiers.
